\documentclass[11pt]{article}
\usepackage{ap-lit}
\title{Unit 3: Longer Fiction and Drama}
\author{Paul Wang}
\date{November 2026}
\docpart{Unit 3 \textperiodcentered\ Study Guide \textperiodcentered\ Longer fiction or drama 15--18\% of MCQ}

\begin{document}
\maketitle

\begin{abstract}
Novels and plays are a smaller share of the multiple choice (usually one passage), but they are the whole of the third essay, the Literary Argument, where you choose the work. This guide covers what changes when a text is long or staged: character over a whole work, including the tragic pattern (section 1); setting as a social world (2); the architecture of long works and plays (3); narration in novels and its absence on stage (4); and irony, wit, symbol, and verse as the language of long works (5). Section 6 is a full plan for Question 3. Each section ends with check questions; the key is on the last page.
\end{abstract}

\scoringline{Check questions: \blank\ / \qtotal \qquad Target: 12+ \qquad My Q3 works: \blank[8em]}

\section{Character across a whole work}
\label{sec:1}

In a short passage you see a character at one moment; in a novel or play you see them change, or fail to. The Literary Argument almost always asks how a character, conflict, or choice contributes to the meaning of the \emph{whole} work, so track arcs:
\begin{itemize}
\item What does the character want at the start, and what stands in the way?
\item What decision or recognition changes them, and what does it cost?
\item Who are their \textbf{foils}, and what do the contrasts reveal?
\item What do \textbf{minor characters} do? They often embody an option the protagonist rejects, or comment on the action.
\end{itemize}

\subsection{The tragic pattern}
\begin{terms}
Tragic hero & A figure of high status and real stature whose downfall comes in part from their own choices. \\
Hamartia & The ``error'' or flaw that contributes to the fall (ambition, pride, indecision); not simply a vice. \\
Peripeteia & The reversal of fortune. \\
Anagnorisis & The recognition: the moment the hero sees the truth of their situation, often too late. \\
Catharsis & The audience's release of pity and fear. \\
\end{terms}

\begin{q}{2}{3.1 Foil}{Contrasting traits that sharpen the protagonist = foil.}
In a novel, a cautious, rule-following brother appears beside an impulsive one, and each scene between them sharpens the impulsive brother's choices. The cautious brother functions mainly as a
\chh{tragic hero}{foil}{narrator}{antagonist}
\end{q}

\begin{q}{4}{3.1 Anagnorisis}{The hero's moment of recognizing the truth = anagnorisis.}
Near the end of a tragedy, a king realizes that the prophecy he trusted was a trap and that his own choices have ruined him. This moment of recognition is called
\chh{catharsis}{peripeteia}{hamartia}{anagnorisis}
\end{q}

\begin{q}{1}{3.1 Minor characters}{Minor characters often embody an alternative the protagonist rejects.}
A minor character who accepts the marriage the heroine refuses most likely serves to
\ch{embody the path the heroine rejects, clarifying what her refusal costs and means}{provide comic relief only}{narrate the novel}{replace the heroine as protagonist}
\end{q}

\section{Setting as a social world}
\label{sec:2}

In long works, setting is often a \textbf{society}: its classes, rules, money, and expectations. Characters are measured against it. A comedy of manners studies the codes of a social class; a novel of the city may set a newcomer against its indifference; a play may confine its characters to one room so that the pressure builds. Watch for \textbf{changes of setting} across a work: moving from country to city, from home to exile, often marks a change in the character.

\begin{q}{3}{3.2 Social setting}{The rules about money and marriage are the pressures the characters push against.}
A novel opens by explaining that in its town, young women without money must marry well or face poverty. This detail of setting chiefly
\ch{establishes the weather and season}{identifies the narrator}{establishes the social and economic pressures that will shape the characters' choices}{reveals the climax of the plot}
\end{q}

\begin{q}{2}{3.2 Change of setting}{A move from one world to another often marks a change in the character.}
A play's first act takes place in a crowded family kitchen; its last takes place in an empty train station. The change most likely marks
\ch{a lower production budget}{the protagonist's movement from family belonging to isolation or independence}{a flashback}{the playwright's preference for public places}
\end{q}

\section{The architecture of long works}
\label{sec:3}

\subsection{Novels}
Novels organize time across chapters. Look for \textbf{parallel plots} and \textbf{subplots} that echo or contrast the main story; for \textbf{framing} (a story told by someone looking back); for \textbf{foreshadowing} and \textbf{motifs} that tie distant chapters together; and for the relationship between the \textbf{opening} and the \textbf{ending}, which often answer each other.

\subsection{Plays}
A play is divided into \textbf{acts} and \textbf{scenes}. Because there is usually no narrator, \textbf{exposition} must come through dialogue, often in an opening conversation. The classic shape rises to a \textbf{climax} and falls to a \textbf{denouement}. In Shakespeare, the climax often comes in Act III, and the last act brings the consequences home. \textbf{Comic relief} (a lighter scene inside a tragedy) sharpens the tragedy around it.

\begin{q}{4}{3.3 Subplot}{A parallel subplot echoes or contrasts the main plot and comments on it.}
In a play, a secondary family's quarrel over an inheritance mirrors the royal family's quarrel over the crown. This subplot chiefly serves to
\ch{give the actors a rest}{introduce the play's narrator}{delay the climax without adding meaning}{echo the main conflict, extending its meaning beyond one family}
\end{q}

\begin{q}{1}{3.3 Exposition in drama}{With no narrator, a play delivers background through dialogue.}
Two servants open a play by discussing their master's recent marriage and his quarrel with his brother. The scene mainly provides
\chh{exposition}{the denouement}{the climax}{an epilogue}
\end{q}

\section{Narration in novels, and its absence on stage}
\label{sec:4}

\subsection{Novels}
The narrators of nineteenth-century novels often know everything and say so: they comment, judge, and generalize, often with \textbf{irony}. A narrator may also move close to a character through \textbf{free indirect discourse}, so the reader shares a character's view while the narrator's framing exposes its limits. Ask always: whose judgment is this sentence carrying?

\subsection{Plays}
Drama has no narrator, so a play reveals inner life through:
\begin{terms}
Soliloquy & A character alone on stage speaks thoughts aloud; we hear what no other character hears. \\
Aside & A brief remark to the audience or to one character, unheard by others on stage. \\
Monologue & A long speech to other characters. \\
Stage directions & The playwright's instructions on setting, action, and manner, often revealing. \\
Dramatic irony & The audience knows what a character does not, as when a character trusts someone we have seen plotting. \\
\end{terms}

\begin{q}{3}{3.4 Soliloquy}{Alone on stage, speaking thoughts aloud = soliloquy.}
Alone on stage, a character debates aloud whether to betray his friend. The speech is a
\chh{dialogue}{monologue to the court}{soliloquy}{stage direction}
\end{q}

\begin{q}{2}{3.4 Dramatic irony}{The audience knows of the poisoned cup; the character does not.}
The audience has watched a villain poison a cup of wine; the hero then raises it to toast his friends. The tension comes chiefly from
\chh{an aside}{dramatic irony}{a soliloquy}{comic relief}
\end{q}

\begin{q}{4}{3.4 Ironic narrator}{A narrator's general ``truth'' that the story shows to be a local prejudice = irony.}
A novel opens with the narrator declaring a ``universal truth'' that, it soon becomes clear, is only the obsession of a few local families. The opening is best described as
\ch{a direct statement of the author's belief}{a soliloquy}{an objective report of facts}{ironic, exposing a social assumption by pretending to endorse it}
\end{q}

\section{The language of long works}
\label{sec:5}

\subsection{Irony, wit, and satire}
\textbf{Verbal irony} (saying one thing and meaning another) is the signature of comic and satiric novels. \textbf{Wit} is quick, pointed verbal cleverness; an \textbf{epigram} is a short, witty, often paradoxical statement. \textbf{Satire} uses irony and exaggeration to criticize individuals or institutions. In a comedy of manners, dialogue reveals character through what people say \emph{and} what they fail to notice.

\subsection{Symbol and motif}
Across a long work, a \textbf{motif} recurs and gathers meaning: blood and sleep in \emph{Macbeth}, the green light in \emph{The Great Gatsby}. A motif's meaning often changes: the same image can promise at the start and mock at the end.

\subsection{Verse and prose in Shakespeare}
Noble characters usually speak \textbf{blank verse} (unrhymed iambic pentameter); servants, clowns, and characters in madness or disguise often speak \textbf{prose}. A shift between them signals a shift in status, control, or state of mind. A \textbf{shared line}, split between two speakers, shows them answering quickly; a short, incomplete line can mark a silence or a shock.

\begin{q}{1}{3.5 Verbal irony}{Praising what is plainly foolish while meaning the opposite = verbal irony.}
A narrator calls a character's endless complaints about her nerves ``her most faithful companions.'' This is
\chh{verbal irony}{an epigram}{dramatic irony}{a soliloquy}
\end{q}

\begin{q}{3}{3.5 Motif}{A recurring image that gathers and changes meaning across a work = motif.}
Images of hands being washed recur in a play, first literally and later in a character's sleepwalking. Such a recurring image is best called
\chh{an aside}{a subplot}{a motif}{exposition}
\end{q}

\begin{q}{2}{3.5 Verse and prose}{A shift from verse to prose often signals a change in control or mental state.}
A noble character who has spoken in blank verse throughout a play begins speaking in broken prose in a late scene. The change most likely signals
\ch{that a different actor has taken the role}{a breakdown in the character's composure or state of mind}{that the scene is a flashback}{that the character has become a servant}
\end{q}

\section{Question 3: the Literary Argument}
\label{sec:6}

\subsection{The prompt}
The prompt opens with a sentence or two about an idea (a character who must choose between two loyalties, a moment of deception, an act of generosity that costs something), asks you to choose a work of fiction in which it appears, and then asks you to analyze how that element \textbf{contributes to an interpretation of the work as a whole}. Some prompts list suggested works; you may choose another of comparable literary merit. Do not write about a film or a work you know only from a summary.

\subsection{Choosing the work}
Choose the work that fits the prompt most fully, not the one you like most. Before the exam, prepare two or three works (one novel and one play is a good pair) and know for each: the central conflicts; three or four key scenes and their order; the ending and its final image; the protagonist's change; two symbols or motifs; and what the work, finally, says.

\subsection{Writing it}
\begin{enumerate}[leftmargin=1.6em]
\item \textbf{Thesis}: name the work and author, the element the prompt names, and the interpretation of the whole work it supports.
\item \textbf{Body}: each paragraph makes a claim, supports it with a specific scene, decision, or image (quotations are welcome but not required), and explains how it builds the interpretation.
\item \textbf{The whole work}: include the ending. An argument that never reaches the end of the book rarely reaches its meaning.
\end{enumerate}

\begin{study}{A sample thesis}
In F.~Scott Fitzgerald's \emph{The Great Gatsby} (1925), Gatsby's careful fabrication of his own past, from the invented inheritance from ``wealthy people in the Middle West'' to the mansion bought so that Daisy would be just across the bay, is not mere fraud but the novel's clearest image of the American belief that a self can be built from desire alone, a belief the novel admires and shows to be fatal.
\end{study}

\begin{trap}
Plot summary is the main reason Q3 essays score 1 or 2 in Row B. A useful test: after every sentence that tells what happens, the next sentence should tell what it means.
\end{trap}

\begin{q}{4}{3.6 Q3 thesis}{Names the work, the prompt's element, and an interpretation of the whole work.}
A prompt asks how a character's act of deception contributes to the meaning of a work as a whole. Which thesis best responds?
\ch{Deception is a common theme in many novels.}{In \emph{Macbeth}, there is a lot of deception.}{\emph{Macbeth} is a tragedy by Shakespeare about a Scottish king.}{In \emph{Macbeth}, the hero's welcoming of Duncan to a castle where he plans his murder shows how ambition corrupts the bond between host and guest, and the play's collapse of order follows from that first betrayal.}
\end{q}

\begin{q}{1}{3.6 Q3 evidence}{Specific scenes and choices, explained in relation to the prompt, are strong Q3 evidence.}
Which is the best kind of evidence for a Literary Argument essay?
\ch{Specific scenes, decisions, and images from the chosen work, each explained in relation to the prompt}{A summary of the author's biography}{Quotations from published critics}{A retelling of the whole plot in order}
\end{q}

\makekey
\end{document}
